% Milestone 1 report, max 2 pages. ACL style files (acl.sty, acl_natbib.bst) are in this folder:
% compile with pdflatex, or upload the whole report/ folder to Overleaf.
\documentclass[11pt]{article}
\usepackage{acl}
\usepackage{times,latexsym,booktabs}
\usepackage[T1]{fontenc}
\usepackage[utf8]{inputenc}

\title{Faithful Explanations for Code-Switched Sentiment Analysis: Milestone 1}
\author{Mahmoud \\ German International University}

\begin{document}
\maketitle

\begin{abstract}
LLM self-explanations can look convincing without reflecting what drove the prediction. We study
whether this holds for Egyptian Arabic--English code-switched sentiment, comparing LLM
self-explanations with Integrated Gradients (IG) on a fine-tuned encoder using erasure-based
faithfulness scores. On EESA, fine-tuned encoders reach 0.88 macro-F1, and in a small CPU-scale study (50 sentences) Integrated Gradients is more faithful than a small LLM's self-explanations, which are no better than random words. Results are preliminary: one run, small models, a single annotator.
\end{abstract}

\section{Introduction and research question}
Self-explanations produced by LLMs are often judged by their plausibility to humans, not by
whether they reflect the evidence the model used \citep{jacovi2020faithfully,turpin2023language}.
Prior work compares extractive self-explanations with human rationales in several text
classification tasks \citep{brandl2024systematic} and tests explanation faithfulness for emotion
detection in Persian \citep{mehrazar2026faithfully}, but code-switched input, where the language of
the explanation is itself a design choice, is unexamined. \textbf{RQ:} How faithful are LLM
self-explanations and gradient-based attributions on code-switched sentiment data, and does
rationale-supervised training improve faithfulness?

\section{Related work}
\citet{brandl2024systematic} find that agreement between LLM rationales and human rationales
depends on text length and task, and that post-hoc attributions often highlight structural tokens.
\citet{mehrazar2026faithfully} study Predict-then-Explain versus Explain-then-Predict prompting in
Persian and find explanations often diverge from faithful reasoning. \citet{turpin2023language} and
\citet{lanham2023measuring} show unfaithful chain-of-thought explanations, and
\citet{jain2019attention} argue that attention weights are not explanations.
\citet{jacovi2020faithfully} argue faithfulness must be evaluated separately from plausibility.
\citet{sundararajan2017axiomatic} introduce Integrated Gradients, and \citet{deyoung2020eraser}
define the erasure-based comprehensiveness and sufficiency scores we use.

\section{Data}
\textbf{EESA} \citep{sherif2024eesa} contains 4,100 Egyptian Arabic--English code-switched YouTube comments with
three sentiment labels, released with official train/dev/test files on GitHub
(\texttt{ASTalaat/\allowbreak EESA-Corpus}). The repository states no licence, so we use it for coursework only and cite the source.
\textbf{Pipeline} (\texttt{src/data.py}): URLs, user mentions and diacritics removed, duplicates dropped inside and across splits
(2,463 / 817 / 816 sentences after cleaning). Labels (train): 1,092 positive, 777 neutral, 594 negative. Mean length 12.6 words;
all sentences mix Arabic and Latin script by our script heuristic. Seed 42.
\textbf{Rationales:} for 50 test sentences (random, seed 42) we marked the words that justify the label.
No second annotator
exists, so no agreement is reported (limitation).

\section{Baselines}
Three types, two models each. The cloud machine used had no GPU, so the LLMs are small (4 CPU cores); larger models need a Colab T4:
\textbf{embeddings} (multilingual-e5-base, BGE-M3, frozen, logistic regression);
\textbf{fine-tuned encoders} (MARBERTv2 \citep{abdulmageed2021arbert}, AraBERTv02-twitter);
\textbf{zero-shot LLMs} (Qwen2.5-1.5B-Instruct, Qwen3-1.7B; 200 test sentences), where the label is the
highest next-token probability among the label words. Metrics: accuracy and macro-F1 on the test set (one run, seed 42).

\begin{table}[h]\centering\small
\begin{tabular}{llcc}\toprule
Type & Model & Acc. & Macro-F1\\\midrule
Embedding & multilingual-e5-base & .824 & .812\\
Embedding & BGE-M3 & .830 & .817\\
Encoder & MARBERTv2 & .887 & .880\\
Encoder & AraBERTv02-twitter & .876 & .871\\
LLM & Qwen2.5-1.5B-Instruct & .625 & .584\\
LLM & Qwen3-1.7B & .710 & .672\\\bottomrule
\end{tabular}\caption{Test results on EESA (LLMs: 200 sentences).}
\end{table}

\section{Limitation experiment: faithfulness}
For each sentence the LLM predicts a label and lists up to five words that influenced it. Using the
LLM's own probability $p$ of the predicted label, comprehensiveness is $p(x)-p(x\setminus R)$ and
sufficiency is $p(x)-p(R)$ \citep{deyoung2020eraser}, where $R$ is the explanation word set; a random set of the same
size is the control. For the fine-tuned encoder we compute IG on the input embeddings
\citep{sundararajan2017axiomatic} (32 steps, all-PAD baseline, token scores summed to words, completeness
error reported) and apply the same erasure test to the top-$k$ words, with $k$ equal to the human
rationale length. Both are also compared to the human rationales (word-level F1).
Table~2 shows results on the 50 sentences (random = same-size random word set). IG on both encoders removes much more
confidence than random words (comprehensiveness .24/.21 vs .08/.09) and the top words alone preserve more
(sufficiency .12/.05 vs .23/.24). The Qwen3-1.7B self-explanations are \emph{not} better than random: erasing
them drops confidence by .28 versus .36 for random words, and keeping only them preserves .34 versus .36.
Agreement with the human rationales is higher for AraBERT-IG (.47 F1) than for the LLM (.41). Caveats: the LLM and the encoders
are different models, so scores are not directly comparable; one run; $n{=}50$; no confidence intervals yet.

\begin{table}[h]\centering\small\setlength{\tabcolsep}{3pt}
\begin{tabular}{lccccc}\toprule
 & Comp.$\uparrow$ & (rand) & Suff.$\downarrow$ & (rand) & F1\\\midrule
Qwen3-1.7B self-expl. & .282 & .363 & .336 & .359 & .410\\
IG MARBERTv2 & .245 & .083 & .116 & .228 & .388\\
IG AraBERT & .213 & .094 & .047 & .236 & .469\\\bottomrule
\end{tabular}\caption{Erasure faithfulness and word-level F1 with the human rationales.}
\end{table}

\section{Method v1 and plan}
Method v1 (rationale-supervised fine-tuning with an attribution-alignment loss on the rationales) is planned, not implemented.
Next: run larger LLMs (7--8B, 4-bit) on a GPU, add confidence intervals and several seeds,
test explanation language (Arabic vs.\ English) as a design choice.

\section*{Limitations and AI use}
No inter-annotator agreement score; small CPU-scale models;
one run. An AI assistant was used for code, experiments, and parts of the report.

\bibliographystyle{acl_natbib}
\begingroup\small\setlength{\bibsep}{1pt}
\begin{thebibliography}{9}
\bibitem[Abdul-Mageed et~al.(2021)]{abdulmageed2021arbert}
Muhammad Abdul-Mageed, AbdelRahim Elmadany, and El~Moatez Billah Nagoudi. 2021.
ARBERT \& MARBERT: Deep bidirectional transformers for Arabic. In \emph{Proc.\ ACL-IJCNLP}, pages 7088--7105.
\bibitem[Brandl and Eberle(2024)]{brandl2024systematic}
Stephanie Brandl and Oliver Eberle. 2024. A systematic comparison between extractive self-explanations and human rationales in text classification. arXiv:2410.03296.
\bibitem[DeYoung et~al.(2020)]{deyoung2020eraser}
Jay DeYoung, Sarthak Jain, Nazneen~Fatema Rajani, Eric Lehman, Caiming Xiong, Richard Socher, and Byron~C. Wallace. 2020. ERASER: A benchmark to evaluate rationalized NLP models. In \emph{Proc.\ ACL}, pages 4443--4458.
\bibitem[Jacovi and Goldberg(2020)]{jacovi2020faithfully}
Alon Jacovi and Yoav Goldberg. 2020. Towards faithfully interpretable NLP systems: How should we define and evaluate faithfulness? In \emph{Proc.\ ACL}, pages 4198--4205.
\bibitem[Jain and Wallace(2019)]{jain2019attention}
Sarthak Jain and Byron~C. Wallace. 2019. Attention is not explanation. arXiv:1902.10186.
\bibitem[Lanham et~al.(2023)]{lanham2023measuring}
Tamera Lanham et~al. 2023. Measuring faithfulness in chain-of-thought reasoning. arXiv:2307.13702.
\bibitem[Mehrazar et~al.(2026)]{mehrazar2026faithfully}
Mobina Mehrazar, Mohammad~Amin Yousefi, Parisa Beygi, and Behnam Bahrak. 2026. Can LLMs faithfully explain themselves in low-resource languages? A case study on emotion detection in Persian. In \emph{Proc.\ LREC}, pages 560--579.
\bibitem[Sherif and Sabty(2024)]{sherif2024eesa}
Ahmed Sherif and Caroline Sabty. 2024. Sentiment analysis for Egyptian Arabic-English code-switched text (EESA corpus), SPECOM 2024; data: github.com/ASTalaat/EESA-Corpus.
\bibitem[Sundararajan et~al.(2017)]{sundararajan2017axiomatic}
Mukund Sundararajan, Ankur Taly, and Qiqi Yan. 2017. Axiomatic attribution for deep networks. arXiv:1703.01365 (ICML 2017).
\bibitem[Turpin et~al.(2023)]{turpin2023language}
Miles Turpin, Julian Michael, Ethan Perez, and Samuel~R. Bowman. 2023. Language models don't always say what they think: Unfaithful explanations in chain-of-thought prompting. arXiv:2305.04388.
\end{thebibliography}
\endgroup
\end{document}
